% ===================== sec_discussion.tex =====================
% NOTE: every sentence in this section interprets a result reported in Sections 5 and 6. Explanations that were not tested are written as hypotheses and are marked as such. No causal wording is used for observational comparisons (claim_audit_report.md, Check 9).
\section{Discussion}\label{sec:discussion}

\subsection{Principal findings}\label{sec:disc_findings}
Four findings emerge from the evaluation of twelve models on four imaging modalities. First, five general-purpose models formed the group with the lowest \SMHR{}, and the data do not support an order within the three leading models. A coarse ordering is therefore the most that this benchmark can offer, and a single best model should not be named from it. Second, each of the three domain-adapted models had a higher \SMHR{} than the best general-purpose model in every modality, by a factor of \PH{1.31} to \PH{1.85}. This result remained under equal class weights, under random changes of the weights, under selection of the reference model on half of the images and under symmetric noise in the judge labels. Third, the level of \SMHR{} depends strongly on the modality, with the highest values in radiography \rad{} and the lowest in dermoscopy, and the composition of the errors differs between modalities. Fourth, the experiments on mitigation and calibration gave results that are specific to the model and, in the case of retrieval, depend on the choice of the baseline.

The second finding is the one that bears most directly on the use of domain-adapted models, and it should be read within its limits. The comparison is with the best of nine general-purpose models, and not with all of them. One domain-adapted model, HuatuoGPT-Vision, had a lower \SMHR{} than between one and three of the general-purpose models in every modality. The three domain-adapted models differ from the general-purpose models in size, architecture and date of release, and the design contains no pair of a base model and the model adapted from it. A difference in \SMHR{} cannot therefore be attributed to the adaptation itself. Studies that did use such pairs found no reliable improvement from medical fine-tuning, with results that depended on the task and on the wording of the prompt \citep{mclaughlin2026knowledge,sabbaghziarani2026confident}, and other studies found that medical models can be less robust than general ones \citep{wu2024halluvqa,gu2024medvh,xu2025sycophancy}. The present result is consistent with these reports. It adds a measurement made with one protocol for four modalities, and it does not explain the difference.

\subsection{Possible explanations of the specialisation result}\label{sec:disc_explain}
Several explanations are compatible with the data, and none of them was tested here. We list them to indicate what the design cannot exclude.
\begin{itemize}
  \item \emph{Style of the response.} The \SMHR{} is a mean over claims. Models that produce short answers contribute fewer claims, and each error then carries more weight in the mean. The two domain-adapted models with the shortest outputs, LLaVA-Med and CheXagent, yielded \PH{21,767} and \PH{24,685} claims and omissions over \Ncalls{} responses each, against \PH{78,744} for Qwen3.5. The comparison of retrieval-augmented prompting in Section~\ref{subsec:res_mitigation} shows that the number of correct claims can move the \SMHR{} without a change in the share of responses with a serious error. The \CHR{} does not depend on the number of claims in the same way, and it ordered the domain-adapted models less unfavourably (CheXagent had one of the lowest values), which is consistent with a part of the difference arising from the length and style of the response.
  \item \emph{Mismatch of the prompts.} The queries are free-text instructions in a common format. A domain-adapted model may have been trained on data with a narrower format, for example captions, reports or multiple-choice questions, and may follow an open instruction less reliably.
  \item \emph{Quantisation and decoding.} All models were run with 4-bit quantisation and with sampling at a low temperature. The effect of quantisation may differ between models, and for LLaVA-Med a decoding defect truncated a large share of the responses (Section~\ref{subsec:res_sri}).
  \item \emph{Scope of training data.} CheXagent was developed for chest radiography, and its values in the other three modalities measure in part the transfer of a model for one modality to others. LLaVA-Med was developed for broader biomedical use, and its \SRI{} values in the other three modalities were similar to those of CheXagent, while those of HuatuoGPT-Vision were lower.
  \item \emph{Interaction with the judge.} The judge is a general-purpose language model, and a general-purpose judge may rate the typical output of a general-purpose model more favourably. The audit could not test this (Section~\ref{sec:val_summary}).
\end{itemize}
A study that settles the question requires base--specialist pairs, full-precision inference, repeated sampling and a judge from a different family, and we return to this in the final section.

\subsection{Implications for the evaluation of medical vision-language models}\label{sec:disc_eval}
\paragraph{Weighting by severity changes the interpretation more than the ranking.} The ranking of the models by \SMHR{} was nearly unchanged when the class weights were set to be equal or perturbed by up to 30\%, and it was identical to the ranking by the MediHall Score, which is a weighted error rate over the same classified claims. On these data, severity weighting therefore does not by itself reorder the models. Its value lies in what it reports: the \SMHR{} states how much of the error comes from the classes that carry the highest potential for harm, and the class composition shows that fabrication and omission account for between \PH{70\%} and \PH{88\%} of the flagged claims in the four modalities. A reader who accepts the order of the classes can interpret the score. A reader who does not can recompute it with other weights, which is possible if the class labels are released with the benchmark.

\paragraph{A rate per claim and a rate per response answer different questions.} The \SMHR{} and the \CHR{} were not correlated across models ($\rho = \PH{0.00}$), and, against the benchmark baseline, retrieval-augmented prompting lowered or left unchanged the \SMHR{} of LLaVA-1.6, HuatuoGPT-Vision and LLaVA-Med while it raised their \CHR{}. A model can add correct statements to its answer and lower the rate per claim while more of its answers contain a serious error. A benchmark that reports only a rate per claim, or only a rate per response, can be improved on the reported measure by changes that do not improve the answers. We recommend that both be reported, together with the number of claims.

\paragraph{Mitigation studies need a matched baseline.} Changing the decoding setting from sampling at a temperature of 0.1 to greedy decoding changed the \SMHR{} of LLaVA-1.6 by \PH{$-0.046$}. This is comparable in size to the apparent effects of the mitigation strategies, and it reversed the sign of the apparent benefit of retrieval for the same model. The only effect of mitigation that held against both baselines, with and without the radiology images, was the reduction in \SMHR{} for HuatuoGPT-Vision under contrastive decoding. The other effects were specific to a model, small or of uncertain sign, and the experiment does not support a general recommendation for either strategy.

\paragraph{Stated confidence is not a usable safety signal for these models.} The accuracy of the responses under the strict label was between \PH{12\%} and \PH{28\%}, while the confidence that the models stated averaged between \PH{52\%} and \PH{91\%}. The label is demanding, since a response is counted as correct only if it contains no hallucination of any class, but the gap is large for every model and configuration, and it came mostly from overconfident bins. The order of the models by calibration error varied little with the decoding setting, and HuatuoGPT-Vision was the worst-calibrated model in every configuration. With one prompt and five models, we regard this as a description of the models examined and not as a property of their weights.

\paragraph{The modality changes what the score means.} The \SMHR{} reflects different mixes of error in different modalities, and the share of the weighted error that can be checked against the reference labels varied between modalities. In histopathology, almost all of the weighted error arose from classes that could be checked only in part, because half of the patches carry a binary tumour label alone. A value of \SMHR{} in one modality is therefore not the same quantity as a value in another, and the \SRI{}, which compares models within a modality, is more interpretable across modalities than the absolute rates.

\subsection{Interpretation for clinical use}\label{sec:disc_clinical}
The metrics describe the disagreement of model statements with the reference labels, weighted by an assumed order of harm. They do not describe the harm to patients, they do not use clinical outcomes, and they do not establish that any model is safe or unsafe in practice. The reference labels are coarse, since they record a diagnosis or a small set of findings and not a complete reading of the image, and a statement that is absent from the labels is counted as unsupported even when it may be correct. The finding that the domain-adapted models had higher rates than the best general-purpose models does not imply that general-purpose models are suitable for clinical use, because the rates of the best group are themselves high (\SMHR{} from \PH{0.529}) and the proportion of responses with a fabrication or an omission was between \PH{72\%} and \PH{81\%} for all twelve models. We regard the results as evidence of the size of the problem in the models examined, and as a method for measuring it.
